% Mục 5 (bản rút gọn). Rút từ sections/method.tex; nguồn từng số ghi ở đó. Nguồn chính: scripts/run_bench.py,
% include/bench_common.hpp, scripts/plot.py, results/<máy>/{calibration.json, campaign.log, *.csv}.
Chương trình đo lock chạy vòng lặp của đề, mỗi thread $I$ vòng. Chương trình đo counter thay critical
section bằng một lần tăng counter, với OW lời gọi \texttt{is\_prime(101)} giữa hai lần tăng. Chương trình
đo barrier gọi barrier $I$ lần liên tiếp, không có việc ở giữa, nên mỗi lần chạy có đúng $I$ episode.

Đồng hồ là \texttt{std::chrono::steady\_clock}. Vùng đo bắt đầu khi mọi thread đã được tạo và cổng
xuất phát chung mở ra, và kết thúc khi lần \texttt{join} cuối trả về. Đề yêu cầu thời gian của cả
chương trình. Báo cáo dùng vùng đo làm số đo chính, vì việc khởi động tiến trình và tạo thread không
thuộc chi phí đồng bộ mà các cài đặt cần so. Runner cũng ghi thời gian cả chương trình, và
Bảng~\ref{tab:cross:endtoend} so hai thời gian này. Khi số thread nhỏ hơn số luồng phần cứng, thời gian
cả chương trình dài hơn vùng đo trung vị 3--4~ms trên hai máy Apple và 8--10~ms trên \PCname{LAP1184}.
Từ số luồng phần cứng trở lên, phần chênh là 16--42~ms. Với lock, thứ tự giữa hai lock cùng cấu hình đổi
ở khoảng 2\,\% số cặp, và chỉ giữa các cặp chênh nhau không quá 5\,\%. Với counter trên
\PCname{LAP1184}, thứ tự đổi ở 6\,\% số cặp, và có cặp chênh nhau tới 36\,\%. Độ phân giải đồng hồ
rất nhỏ so với một lần chạy, có độ dài trung vị 0,15--0,29~s. Trên hai máy Apple đồng hồ đếm ở 24~MHz,
tức 41,7~ns một nhịp (trung vị đo được 41,0~ns). Trên \PCname{LAP1184} mỗi lần đọc đồng hồ tốn 20~ns.
% nguồn: report/figures/cross/tables.tex bảng endtoend (cột wall_s của results/<máy>/*.csv); sleep_latency.csv
% dòng clock; scripts/report_numbers/corr_stats.py.

Mỗi lần chạy in một dòng CSV gồm tham số, thời gian, load average 1 phút và kết quả của một phép kiểm.
Phép kiểm buộc mọi phần việc phải thật sự chạy: \texttt{is\_prime} không inline, \texttt{N} đi qua
\texttt{pc::opaque()} ở mỗi lời gọi, và tổng kết quả của mỗi thread phải bằng số lời gọi. Trên
\PCname{apple-m5} ở $N=1009$, thời gian một lời gọi $t$ nằm trong 504--527~ns ở mọi OW từ 0 tới 15 (trên
\PCname{LAP1184} là 1\,600--1\,640~ns), nên phần việc không bị xoá. Trên
cả ba máy, mọi lần chạy hoàn thành đều qua phép kiểm kết quả.

Phần lock quét 12 lock qua ba lượt. Một lượt dùng $N=1009$ với OW bằng 0, 1, 3, 7, 15. Hai lượt còn lại
dùng $N=11$ và $N=10\,007$ với OW bằng 0 và 3. Ba giá trị cho critical section khoảng 5~ns, 0,5~\textmu s
và 5~\textmu s. Barrier quét 7 cài đặt. Counter quét 19 cài đặt với OW bằng 0, 1, 8. Số thread trên
\PCname{apple-m5} là 1, 2, 3, 4, 6, 8, 10, 12, 16, 20, trong đó 4 là số P-core, 10 là mọi core và 12--20
là quá tải. Hai máy còn lại thêm 24 thread.

Với lock, số vòng mỗi thread là $I = W_N/\max(p,\,1+\mathrm{OW})$, với $W_N$ hiệu chỉnh một lần cho mỗi
$N$ để một lần chạy của \PCname{ttas} với OW $=0$ dài khoảng 0,19~s. Với barrier và counter, mỗi cấu hình
có $I$ riêng suy từ một lần chạy thử, nhắm 0,25~s và 0,15~s. $I$ cố định cho cả đợt đo, và hình quy thời
gian về cùng số vòng lặp. Thời hạn mỗi lần chạy là 30~s cho lock và 10~s cho barrier và counter. Riêng
\PCname{sr} và \PCname{sr_bo_spin} khi vượt số core được 30~s, vì một episode có thể tốn cả một lượt lập
lịch. Lock MCS quá tải được đo với số vòng nhỏ hơn 100 lần ngay từ đầu. Lần vượt thời hạn được ghi lại, cấu hình đó
được đo tiếp với số vòng nhỏ hơn 100 lần (không dưới 10), và trung vị chỉ dùng các lần chạy có cùng số
vòng. Mọi lần vượt thời hạn trên ba máy là MCS hoặc \PCname{lock_mcs} quá tải, và không cấu hình nào có
trung vị là một lần vượt thời hạn. Vài ô chỉ còn 4--6 giá trị: \PCname{lock_mcs} quá tải của counter
trên cả ba máy, và MCS quá tải với $N = 11$ trên \PCname{LAP1184}.

Mỗi cấu hình chạy 7 lần. Báo cáo dùng trung vị, không dùng trung bình, vì phân bố lệch phải. Trên
\PCname{apple-m5}, ở 33\,\% cấu hình lock lần chậm nhất dài hơn 1,5 lần trung vị, với barrier là 49\,\%.
Trên hai máy còn lại, tỉ lệ này ở lock là 3\,\% và 7\,\%. Lượt lặp là vòng ngoài.
Trong mỗi lượt, mọi cấu hình chạy theo thứ tự xáo trộn với seed cố định, nên tải nền trôi chậm không
thiên vị cài đặt nào. Mỗi thời điểm chỉ có một tiến trình đo.

Phần lock được so với cận dưới lý tưởng
\[ T_{\min} = \max(p,\ 1+\mathrm{OW})\cdot I\cdot t. \]
Hai vế trong $\max$ là hai cận dưới riêng. Các critical section loại trừ nhau, nên tổng thời gian của
chúng, $p\,I\,t$, là một cận dưới. Mỗi thread phải tự gọi \texttt{is\_prime} $(1+\mathrm{OW})\,I$ lần,
nên $(1+\mathrm{OW})\,I\,t$ là cận dưới thứ hai. Từ $p = 1+\mathrm{OW}$ vế đầu lớn hơn, và thời gian
đo chỉ chạm cận khi lock bận liên tục. $t$ là trung vị lần chạy 1 thread của lock nhanh nhất ở cùng $N$
và OW, chia cho $I(1+\mathrm{OW})$.
\begin{itemize}[nosep]
  \item \emph{Độ chậm} là thời gian đo chia cho $T_{\min}$. Độ chậm 1 nghĩa là không có chi phí nào
    ngoài phần việc. Phần vượt 1 gồm chi phí của lock và cả các chi phí không do lock, nêu ở mục Các
    yếu tố đe dọa độ tin cậy.
  \item \emph{Độ công bằng} là thời gian của thread xong sớm nhất chia cho thời gian của thread xong
    muộn nhất trong một lần chạy. Mọi thread làm cùng số vòng, nên 1 là hoàn toàn đều.
\end{itemize}

Mỗi điểm trên hình là trung vị của 7 lần lặp, và thanh sai số đi từ min tới max. Nếu 7 lần lặp độc
lập và cùng phân bố, khoảng từ min tới max chứa trung vị của phân bố với xác suất $1 - 2 \cdot 2^{-7}
\approx 98{,}4\,\%$. Ở các ô chỉ còn 4--6 lần lặp, xác suất này là 87,5--96,9\,\%. Vạch dọc liền đánh dấu số core, hoặc số luồng phần cứng khi máy
có SMT, còn vạch dọc chấm đánh dấu số P-core, hoặc số core do máy ảo báo. Vòng tròn lớn quanh một điểm
cho biết cấu hình đó được đo với số vòng giảm 100 lần hoặc hơn rồi quy đổi tuyến tính.
